%% CV for AskCody - Full-Stack Software Engineer (Dansk)
%% Compiled with: lualatex

\documentclass[11pt,a4paper,sans]{moderncv}
\moderncvstyle{banking}
\moderncvcolor{blue}

\renewcommand*{\firstnamestyle}[1]{{\fontsize{34}{36}\bfseries\upshape\color{color1}#1}}
\renewcommand*{\lastnamestyle}[1]{{\fontsize{34}{36}\bfseries\upshape\color{color1}#1}}
\renewcommand*{\sectionstyle}[1]{{\sectionfont\color{color1}#1}}

\usepackage[utf8]{inputenc}
\usepackage{hyperref}
\hypersetup{
    colorlinks=true,
    linkcolor=blue,
    filecolor=magenta,
    urlcolor=blue,
    pdftitle={Jacob El-Omar - CV},
    pdfpagemode=FullScreen,
}
\usepackage[scale=0.80]{geometry}
\usepackage{import}

\name{Jacob}{El-Omar}
\address{Aalborg, Danmark}{}{}
\phone[mobile]{+45 53 35 27 82}
\email{elomarjc@gmail.com}
\extrainfo{\href{https://linkedin.com/in/jacob-el-omar}{LinkedIn} \textbullet{} \href{https://github.com/elomarjc}{GitHub}}

\begin{document}

\makecvtitle

% ============================================================
%     PROFIL
% ============================================================

\vspace{6pt}
\small{Nyuddannet civilingeniør (cand.polyt.) i elektroniske systemer fra Aalborg Universitet med stærke fuldstændige udviklingskompetencer inden for C\#, .NET Core og relationelle databaser. Erfaren i at opbygge cross-platform løsninger med Flutter, integrere REST-API'er og anvende moderne AI-first softwareudvikling (prompts, validering og parprogrammering med kodningsagenter). En nysgerrig, PBL-trænet holdspiller, der er klar to levere værdi på tværs af hele jeres R\&D-infrastruktur.}

% ============================================================
%     FAGLIGE KOMPETENCER
% ============================================================

\section{Faglige Kompetencer}
\vspace{1pt}
\begin{itemize}

\item \textbf{C\# \& .NET Udvikling}: C\#, .NET Core, API-design, objektorienteret programmering, objekt-relationel mapping (Entity Framework).

\item \textbf{Full-Stack \& Frontend}: Flutter, Dart, JavaScript, TypeScript, Node.js, HTML/CSS samt kendskab til React og GraphQL-kald.

\item \textbf{AI-First Software Engineering}: Erfaren i at udnytte AI-kodningsagenter (f.eks. Antigravity) aktivt i den daglige udvikling til hurtig prototypering, testdækning og debugging.

\item \textbf{Databaser \& DevOps}: PostgreSQL, SQLite, MySQL, Supabase, samt versionsstyring i Git og CI/CD processer (GitHub Actions).

\item \textbf{Projektstyring \& Kommunikation}: Trænet i agilt teamsamarbejde (Scrum) under AAU's problembaserede læringsmodel (PBL).

\end{itemize}

% ============================================================
%     ERHVERVSERFARING
% ============================================================

\section{Erhvervserfaring}
\vspace{3pt}
\begin{itemize}

\item{\cventry{Jan 2026--Nu}{IT- \& Webudvikler (Frivillig)}{Dawah Danmark}{Aalborg, Danmark}{}{\vspace{1pt}
\begin{itemize}
    \item Vedligeholdt, moderniserede og overvågede organisationens WordPress-platform.
    \item Optimerede databasearkitekturen for at øge load-hastighed og brugertilfredshed.
\end{itemize}}}

\vspace{3pt}

\item{\cventry{Sep 2024--Jun 2025}{AI/ML Engineer (Praktik)}{Ai Health Highway}{Remote}{}{\vspace{1pt}
\begin{itemize}
    \item Udviklede signalfiltreringsalgoritmer i Python til stetoskoper.
    \item Samarbejdede på tværs af globale teams om at validere og teste softwareløsninger.
    \item Sikrede softwareoverensstemmelse med medicinske kvalitetsstandarder via grundig dokumentation.
\end{itemize}}}

\vspace{3pt}

\item{\cventry{Jun 2023--Aug 2024}{Lager- og Logistikmedarbejder}{Lyn Express}{Aalborg, Danmark}{}{\vspace{1pt}
\begin{itemize}
    \item Håndterede varelogistik og lagerstyringssystemer i et højeffektivt arbejdsmiljø.
\end{itemize}}}

\end{itemize}

% ============================================================
%     UDDANNELSE
% ============================================================

\section{Uddannelse}
\vspace{1pt}
\begin{itemize}

\item{\cventry{Sep 2023--Jun 2025}{Civilingeniør, cand.polyt. (Elektroniske Systemer)}{Aalborg Universitet}{Aalborg, Denmark}{}{\vspace{1pt}
Speciale: ``Adaptive Noise Cancellation For Electronic Stethoscopes.'' Udvikling og test af filtreringsarkitekturer til realtidsplatforme.
}}

\vspace{3pt}

\item{\cventry{Sep 2020--Aug 2023}{Bachelor i teknisk videnskab (Elektronik og IT)}{Aalborg University}{Aalborg, Danmark}{}{\vspace{1pt}
Specialisering i reguleringsteknik. Bachelorprojektet omhandlede autonom regulering af traverskran, og vi designede indlejret faldsikringssoftware på PSoC i C.
}}

\end{itemize}

% ============================================================
%     UDVALGTE PROJEKTER
% ============================================================

\section{Udvalgte Projekter}
\vspace{1pt}
\begin{itemize}

\item{\textbf{Grow a Fish (Fritidsprojekt)}: Egenudviklet spil i Flutter/Dart. Oprettede en Supabase PostgreSQL-database til håndtering af multiplayer-data via REST-API'er og realtidssynkronisering.}

\item{\textbf{Adaptive Noise Cancellation (Speciale, Universitetsprojekt)}: Udviklede en signalbehandlings- og filtreringspipeline i Python til medicinske IoT-stetoskoper. Implementerede LMS/RLS-algoritmer til støjreduktion i realtid.}

\item{\textbf{Embedded Faldsikringssystem (Universitetsprojekt, 4. sem.)}: Designede og programmerede et realtids fall-detection system på en Cypress PSoC mikrocontroller i C. Etablerede sensorintegration via I2C.}

\end{itemize}

% ============================================================
%     SPROG
% ============================================================

\section{Sprog}
\vspace{1pt}
\begin{itemize}
\item Dansk (modersmål), Engelsk (flydende).
\end{itemize}

% ============================================================
%     REFERENCER
% ============================================================

\section{References}
\vspace{1pt}
\begin{itemize}
\item{Afleveres efter aftale.}
\end{itemize}

\end{document}
